% concatenated from 241100739.tex
\documentclass[11pt]{amsart}
\usepackage[margin=2.4cm]{geometry}
\usepackage{amsmath}
\usepackage{amsfonts}
\usepackage{amssymb}
\usepackage{xcolor}
\renewenvironment{proof}{{\bfseries Proof.}}{\qed}
\numberwithin{equation}{section} 
%\theoremstyle{theorem}
\newtheorem{theorem}{Theorem}[section] 
\newtheorem{pro}[theorem]{Proposition} 
\newtheorem{cor}[theorem]{Corollary} 
\newtheorem{lemma}[theorem]{Lemma} 
\newtheorem{thm}[theorem]{Theorem}
\theoremstyle{definition}
\newtheorem{defn}[theorem]{Definition} 
\newtheorem{rk}[theorem]{Remark} 
\newcommand{\gt}{\tilde{\gamma}}
\newcommand{\G}{$\Gamma$ }
\newcommand{\lm}{\lambda}
\newcommand{\n}{\newline}
\newcommand{\bs}{\boldsymbol}
\newcommand{\tb}{\textbf}
\newcommand{\be}{$\beta$}
\newcommand{\la}{$_{\lambda}$ }
\newcommand{\Si}{$\Sigma$}
\newcommand{\ap}{\alpha_{p}}
\newcommand{\imp}{\Rightarrow}
\newcommand{\z}{\mathbb{Z}}
\newcommand{\C}{\mathbb{C}}
\newcommand{\R}{\mathbb{R}}
\newcommand{\ps}{S'}
\newcommand{\g}{\gamma_{p}}
\newcommand{\m}{\gamma^{m}_{p}}
\newcommand{\pc}{$\mathcal C'$}
\newcommand{\ce}{$ \mathcal{C} $}
\newcommand{\zp}{$ \mathbb{Z}_{2}*\mathbb{Z}_{p} $}
\newcommand{\io}{\iota }
\newcommand{\secref}[1]{Section~\ref{#1}}
\newcommand{\thmref}[1]{Theorem~\ref{#1}}
\newcommand{\lemref}[1]{Lemma~\ref{#1}}
\newcommand{\remref}[1]{Remark~\ref{#1}}
\newcommand{\propref}[1]{Proposition~\ref{#1}}
\newcommand{\corref}[1]{Corollary~\ref{#1}}
\newcommand{\eqnref}[1]{~{\textrm(\ref{#1})}}
\newcommand{\tabref}[1]{Table~(\ref{#1})}
\newcommand{\defref}[1]{Definition~(\ref{#1})}
\newcommand{\bl}{\textcolor{blue}}
\newcommand{\br}{\textcolor{red}}
\numberwithin{equation}{section} 
\def \h{{\bf H}} 
\renewcommand{\baselinestretch}{1.2}
\usepackage[backref]{hyperref}

\begin{document}
	\title[Spherical growth of reciprocal classes in the Hecke Groups]{ Spherical growth of reciprocal classes in the Hecke Groups }
	\author{Debattam Das  \and Krishnendu Gongopadhyay}
	\address{Indian Institute of Science Education and Research (IISER) Mohali,
		Knowledge City,  Sector 81, S.A.S. Nagar 140306, Punjab, India}
	\email{debattam123@gmail.com}
    	\address{Indian Institute of Science Education and Research (IISER) Mohali,
		Knowledge City,  Sector 81, S.A.S. Nagar 140306, Punjab, India}
	\email{ krishnendu@iisermohali.ac.in}
	\subjclass[2020]{Primary 20H10; Secondary 11F06, 05E16, 37C35}
	\keywords{ Hecke groups, reversible elements, reciprocal geodesics, growth rate}
	
	
	\begin{abstract}
		Let $\Gamma_p$ denote the Hecke group where $p=2r$, $r>0$. Let $\mathcal{N}_l$ denote the set of conjugacy classes of reciprocal elements of word length $l$ in $\Gamma_p$. We prove that for $l \to \infty$, 
		$$|\mathcal{N}_l| = \mathcal{O}\left(\left\lfloor \tfrac{l+1}{2} \right\rfloor^{s-1} \rho^{\left\lfloor \tfrac{l+1}{2} \right\rfloor} \right),
		$$
		where $\mathcal O$ is the `big O',  $\rho \in [\sqrt{2}, 2]$ is the unique positive real root of
		$$
		p(x) = x^{r+1} - 2\sum_{j=1}^{r-1} x^{r-j} - 1,
		$$
		and $s$ is the maximal multiplicity among the roots of $p(x)$. 
        
        Our method relies on the free product structure of the Hecke group $\Gamma_p$, a combinatorial counting function, and recurrence relations derived from cyclically reduced representatives. We also derive that the growth rate of the primitive reciprocal classes of word length $l$ is in agreement with that of $\mathcal{N}_l$. This work generalizes previous results for odd $p$ and provides an explicit asymptotic bound for all Hecke groups.
	\end{abstract}
	
	\date{\today}
	
	\maketitle	
	
	\section{Introduction}\label{1}

    A non-trivial element $g$ in a group $\Gamma$ is called \emph{reversible} (or \emph{real}) if there exists $h \in \Gamma$ such that $hgh^{-1} = g^{-1}$. If $h$ is an involution (i.e., $h^2 = 1$), then $g$ is said to be \emph{strongly reversible} (or \emph{strongly real}). Every strongly reversible element, except involutions themselves, can be written as a product of two involutions. For various perspectives on reversibility, see \cite{BR}, \cite{LR}, \cite{Sa}, \cite{OS}.

In Fuchsian groups, a torsion-free element is reversible if and only if it is strongly reversible; see \cite[Theorem 1.1]{DG}. Moreover, if $hgh^{-1} = g^{-1}$ for some $g \in \Gamma$, then $h$ must be an involution. In the modular group $\mathrm{PSL}(2, \mathbb{Z})$, such elements are referred to as \emph{reciprocal} (following Sarnak \cite{Sa}), and we adopt this terminology. A conjugacy class that contains a reciprocal element is called a \emph{reciprocal class}.

If $g$ is a hyperbolic reciprocal element, then the axis of $g$ projects to a closed geodesic on $\mathbb{H}^2/\Gamma$. A reciprocal element implies that this geodesic passes through an even-order cone point, such geodesics are called \emph{reciprocal geodesics}. Thus, reciprocal classes correspond bijectively to the equivalence classes of reciprocal geodesics on $\mathbb{H}^2/\Gamma$.

An involution $h$ conjugating $g$ to $g^{-1}$ is called a \emph{reciprocator} of $g$. Reciprocators are unique up to elements in the centralizer of $g$; see \cite{BR}. If a reciprocal element $g$ is a hyperbolic with a reciprocator $h$,  then the subgroup generated  by $g$ and $h$ is an infinite dihedral subgroup of $\Gamma$.  Hence, classifying the reciprocal geodesics amounts to classifying such dihedral subgroups up to conjugacy. Each equivalence class of reciprocal geodesics corresponds to the conjugacy class of a reciprocal element, and thus to a dihedral subgroup generated by $g$ and a reciprocator.


In recent years, the asymptotic growth and distribution of reciprocal geodesics have been studied in a variety of settings,  in the modular group \cite{Sa}, \cite{BK},  on hyperbolic surfaces \cite{ES}, and more generally on negatively curved orbifolds \cite{PP}.     Sarnak investigated reciprocal classes in the modular group $\mathrm{PSL}(2, \mathbb{Z})$, obtaining asymptotic growth and equidistribution results; see \cite{Sa}. Bourgain and Kontorovich further studied the growth rate of \emph{low-lying} reciprocal geodesics on the modular surface \cite{BK}. Erlandsson and Souto extended these results to arbitrary two-dimensional hyperbolic orbifolds; see \cite{ES}. In all these works, the hyperbolic length of geodesics serves as the measure for asymptotic growth. 

Parkkonen and Paulin \cite{PP} have recently generalized the results of Sarnak and Erlandsson--Souto to strongly reversible periodic orbits of the geodesic flow on pinched negatively curved Riemannian orbifolds. Their approach, however, is based on the thermodynamic formalism, marking a significant departure from earlier methods.

	
	Another approach, based on the free product structure of the modular group, was used by Basmajian and Suzzi Valli.  They derived asymptotic growth rates for reciprocal classes using the lengths of the associated combinatorial words; see \cite{BV}.  Following this approach, Marmolejo, in \cite{Ma},  obtained the spherical growth for the reciprocal clasess in a Hecke group $\Gamma_p$ for $p$ odd. Here, by {\em spherical growth} we mean the asymptotic count of the number of distinct reciprocal classes represented by words of exactly a given length. 
    
    
    
    As follows from \cite{DG}, the combinatorics of the reciprocal words are relatively simpler for $p$ odd. This is due to the presence of a unique involution up to conjugacy that conjugates a reciprocal element $g$ to $g^{-1}$. For $p$ even, this is not the case as there are two choices for such an involution, up to conjugacy. Accordingly, it requires further analysis to resolve the complication when $p$ is even. 
	
	In this paper, we have generalized Marmolejo's result for a Hecke group $\Gamma_p$, for $p$ arbitrary. In particular, we obtain an asymptotic estimate for the size of the reciprocal classes in $\Gamma_p$ when $p$ is even.  
	
	Recall that the Hecke group $ \Gamma_{p} $ is a Fuchsian group which is isomorphic to the free product  $ \z_{2} \ast \z_{p}$, and it is generated by 
	the maps
	$$\iota: z\mapsto-\dfrac{1}{z}~~~~~ \\ \hbox{ and } \\ ~ ~~~~ \alpha_p: z\mapsto z+ \lambda_{p}, ~ p \geq 3, \hbox{ where}~ \lambda_{p}=2\cos\pi/p. $$  
	In particular, $\Gamma_3$ is the modular group. Note that the map $ \iota\ap :z \mapsto -\dfrac{1}{z+\lambda_{p}} $ has order  $ p $. We denote $ \g=\iota \alpha_p$. So we can write the presentation of the Hecke group as $\langle~\io,\g~|~\io^2=\g^p=1 ~\rangle$. We will fix this presentation for the Hecke group throughout the paper. When $p=2r$, $\g^r$ is an involution. We denote $\tilde \gamma=\g^r$. 
	
	Since the spherical growth of a reciprocal class for odd $p$ follows from the work of Marmolejo~\cite{Ma}, without loss of generality, we may assume throughout this paper that $p$ is even. Specifically, we take $p = 2r$ with $r > 0$. We prove the following result.
	
	\begin{theorem} \label{mthm}
		Let $\mathcal{N}_l$ denote the set of conjugacy classes of reciprocal elements of word length $l$. Then, as $l \to \infty$,
		$$|\mathcal{N}_l| = \mathcal{O}\left(\left\lfloor \tfrac{l+1}{2} \right\rfloor^{s-1} \rho^{\left\lfloor \tfrac{l+1}{2} \right\rfloor} \right),
		$$
		where $\rho \in [\sqrt{2}, 2]$ is the unique positive real root of
		$$
		p(x) = x^{r+1} - 2\sum_{j=1}^{r-1} x^{r-j} - 1,
		$$
		and $s$ is the maximal multiplicity among the roots of $p(x)$.
	\end{theorem}
    From the above theorem, we have the following corollary.
    \begin{cor}\label{cor1}
	Let $\mathcal{N}^{p}_l$ denote the set of conjugacy classes of primitive reciprocal elements of word length $l$. Then, as $l \to \infty$,
$$|\mathcal{N}^{p}_l| = \mathcal{O}\left(\left\lfloor \tfrac{l+1}{2} \right\rfloor^{s-1} \rho^{\left\lfloor \tfrac{l+1}{2} \right\rfloor} \right),
$$
where $\rho \in [\sqrt{2}, 2]$ is the unique positive real root of
$$
p(x) = x^{r+1} - 2\sum_{j=1}^{r-1} x^{r-j} - 1,
$$
and $s$ is the maximal multiplicity among the roots of $p(x)$.

		\end{cor}
	In addition to the above asymptotic bound, we provide explicit combinatorial formulas for the number of conjugacy classes in each reciprocal type.In addition to the global asymptotic bound, we provide explicit combinatorial formulas for the number of conjugacy classes in each reciprocal type, see \thmref{count1}. We have used Marmolejo's counting function $\Psi^r_n(x)$, cf.~\cite{Ma}, which counts the number of integer solutions to the equation $\sum_{i=1}^n p_i = x$, with $0 < p_i \leq r$. This extends Marmolejo's count to reciprocal classes in $\Gamma_p$ for all values of $p$.
	
	To prove the results, we have used the combinatorial structure of $\Gamma_p \approx \mathbb Z_2 \ast \mathbb Z_p$ to determine the spherical growth.   
	Each reciprocal class contains some cyclically reduced words, which can be expressed in terms of the generators $\iota$ and $\g$. Let $\gt = \g^r$. Then a reciprocal element in $\Gamma_p$ is conjugate to one of the following forms: $\iota k \gt k^{-1}$, $\gt k \gt k^{-1}$, or $\iota k \gt k^{-1}$, for some word $k$. We refer to a word of the form $\iota \g ^{k_1}\io \g^{k_2}\dots \io \g^{k_n}$, that is, starting with $\iota$ and ending with $\g$, as an \emph{$\iota$-$\g$ word}. Using these conjugacy representatives, we derive a summation formula for their word lengths, which in turn yields a recurrence relation associated with each reciprocal class. Finding a closed-form solution for this recurrence relation provides the desired estimate.

 
	
	Now, we briefly outline the structure of this paper. After recalling some preliminary notions, we note preliminary observations about the counting function in \secref{prel}. We estimate counts for the reciprocal classes in \secref{3}. Finally, we prove \thmref{mthm} and \corref{cor1} in \secref{4}.  
	
	
	\medskip
	\tb{Notations:} We have used the following notations throughout the paper.
	\begin{enumerate}
		\item For some real number $x,$ $\lfloor x\rfloor$ is largest integer less than equal to $x$ and $\lceil x\rceil$ is the smallest integer greater than equal to $x$. 
		\item For some real-valued functions $f,g$ such that they are defined on the  same domain $D\subset \mathbb{R}$, we  say \[f(x)=\mathcal{O}(g(x))~~\text{ as }~x\rightarrow\infty \] if there exists $x_0\in D$ and $M>0$ such that \[|f(x)|\leq M|g(x)|~~\text{for}~~x\geq x_0.\] 
	\end{enumerate}
	
	\section{Preliminaries}\label{prel}
	%In \secref{3}, we will see the absolute number of conjugacy classes of different type of reciprocal elements
	
	\subsection{Reciprocal elements in Hecke groups}In \cite{DG}, different types of reciprocal elements for Hecke groups were defined. Here the following types are:
	\begin{defn}
		The reciprocal word is said to be \textit{symmetric} if it can be seen as a product of those elements which are only conjugate to $ \io $.
	\end{defn}
	\begin{defn}
		The reciprocal word is said to be \textit{$ p $-reciprocal} if it can be seen as a product of those elements which are only conjugate to $ \gt $.
	\end{defn}
	\begin{defn}
		The reciprocal word is said to be \textit{symmetric $ p $-reciprocal} if it can be seen as a product of conjugate of $\io$ and conjugate of $ \gt $ respectively.
	\end{defn}
	\begin{rk}
		If $p$ is an odd integer, all reciprocal elements are symmetric. Other types appear in the even case.
	\end{rk}
	\subsection{Word Length} In this section, we will define the word length of the elements of the free products and conjugacy classes.
	\begin{defn}\label{word}
		The\textit{ length} of a word is the number of the generating elements of the group presented in its reduced presentation.
	\end{defn}
	In the group $\Gamma_{p,q} =\langle a, b ~ | ~ a^{p},b^{q} \rangle  $ where $p,q$ are both odd, the length of the reduced word $ a^{j_{1}}b^{k_{1}}a^{j_{2}}b^{k_{2}}\dots a^{j_{n}}b^{k_{n}} $ where $-p/2\leq j_{i}\leq p/2, -q/2\leq k_{i}\leq q/2 $, is $ \sum_{i=1}^{n}|{j_{i}}|+|{k_{i}}| $.
	\begin{defn}
		The \emph{length} of the conjugacy class of a word is the minimum length of a cyclically reduced word present in that conjugacy class. 
		
		If a conjugacy class contains cyclically reduced word with minimal length $l$, we loosely call it a \emph{conjugacy class of length $l$}. 
	\end{defn}
	Our first objective is to calculate the number of conjugacy classes and the number of different reciprocal classes with the same word length  $l.$ To deduce such a number, we have used the following function:
	
	
	
	\subsection{A counting function: } Let the function $\Psi_{n}$ be defined by the number of solutions of the equation, \begin{equation}\label{2.1}
		\sum_{i=1}^{n} p_{i}=x, 
	\end{equation}
	where, $p_{i}$, $i=1, \ldots, n$,  and $x$ are positive integers. It is easy to see that $\Psi_{n}(x) ={x-1 \choose n-1}$. 
	Let $\Psi_{n}^{r}(x)$ denote the number of solutions of \eqnref{2.1} with constraints $p_i \leq r$, $i=1, \ldots, n$,  for some $r<x$. We have the following Proposition.
	
	\begin{pro}
		
		
		Let \( \Psi^r_n(x) \) denote the number of positive integral solutions of the equation, \begin{equation}\label{2.0}
			\sum_{i=1}^{n} p_{i}=x, 
		\end{equation}
		where, $p_{i}$'s, and $x$ are positive integers, and each $p_{i}$  satisfying \( 1 \leq p_i \leq r \). Then,
		\[
		\Psi^r_n(x) = \sum_{k=0}^{\left\lfloor \frac{x - n}{r} \right\rfloor} (-1)^k \binom{n}{k} \binom{x - rk - 1}{n - 1}
		\]
		for \( x \ge n \). If \( x < n \), then \( \Psi^r_n(x) = 0 \).
	\end{pro}
	
	\begin{proof}
		We consider
		in the Equation \eqnref{2.0}  each \( p_i \in [0, r] \). 
		
		Let \( S \) be the set of all solutions with \( p_i \ge 1\), and let \( A_i \) be the set of solutions where \( p_i \ge r+1 \), i.e., where the upper bound is violated in the \( i \)-th coordinate. Let $S^{\prime}$ be the set of all require solutions. Then we have, 
		\[  S^{\prime} =S-(\bigcup_{i}A_i)  \]
		Then the number of require solutions is
		\[
		|S'| = |S| - \sum |A_i| + \sum_{i<j} |A_i \cap A_j| - \cdots + (-1)^k \sum_{i_1 < \dots < i_k} |A_{i_1} \cap \cdots \cap A_{i_k}| + \cdots,
		\]
		by the inclusion-exclusion principle.
		
		We now count the number of solutions with \( k \) variables having the lower bound $r$. There are \( \binom{n}{k} \) ways to choose the violating positions of $p_i$'s, and in each, the sum is \( rk \). So we are solving:
		\[
		\sum p_i = x - rk \quad \text{(with no upper bound on \( p_i \))}
		\]
		which has \( \binom{x - rk - 1}{n - 1} \) solutions.
		
		Hence, the total number is:
		\[
		\Psi^r_n(x) = |S^{\prime}|=\sum_{k=0}^{\left\lfloor \frac{x - n}{r} \right\rfloor} (-1)^k \binom{n}{k} \binom{x - rk - 1}{n - 1}.
		\]
		This completes the proof. 
	\end{proof}
	
	%In this section, the growth of the number of conjugacy classes of hyperbolic words in the Hecke group $\Gamma_{p}$.  %This section is inspired by the work  \cite{BV} where counting of the conjugacy classes in the modular group has been investigated. 
	%The actual value of $\Psi_{n}^{r}(t)$ is not known in general as $\Psi_{n}(t)$. But for some specific cases, it is known. \\
	\begin{lemma}\label{3.1}
		Let $r, x$ be given natural numbers. Then the total number of the solutions $\{n; k_1, \ldots, k_n\}$, $n>0, {k}_i \in \mathbb Z$, $-r<k_{i}\leq r$ for all $k_{i}$ and $q$ be the number of $k_{i}$'s which are equal to $r$ of the equation $\sum_{i=1}^{n}|k_{i}|+n=x$ is given by the following sum after considering all possibilities for $ n$ and $k_{i}$'s, 
		\[\sum_{q=0}^{\lceil\frac{x}{r+1}\rceil-1}\sum_{n=\lceil \frac{x-q}{r}\rceil}^{\lfloor \frac{x-(r-1)q}{2}\rfloor}\Psi_{n}^{r-1}(x-n-rq)2^{n-q} {n\choose q}.\]
	\end{lemma}
	\begin{proof}
		Since some of the $k_{i}$'s are equal to $r$, we have two choices depending upon their sign for those $k_{i}$'s such that $k_{i}\neq r$. So, the number of the solutions for the equation, for fixed $n$ and a fixed $q$ is:  
		\begin{equation}
			\Psi_{n}^{r-1}(x-n-qr)2^{n-q}.
		\end{equation}
		Now, $n$ will be maximum when all $k_{i} \neq r$ have absolute values $1$. Therefore,
		\[n-q+n+rq\leq x\]
		\[\imp 2n+(r-1)q\leq x\]
		\[\imp n\leq \frac{x-(r-1)q}{2}.\]
		Similarly, $n$ will reach the lower bound if the absolute value of each $k_{i} \neq r$ is $r-1$. Therefore,
		\[(n-q)(r-1)+rq+n\geq x\]
		\[\imp nr+q\geq x \]
		\[\imp n\geq \frac{x-q}{r}.\] So,  the number of all the solutions of the equation for fixed $q,$
		\[\sum_{n=\lceil \frac{x-q}{r}\rceil}^{\lfloor \frac{x-(r-1)q}{2}\rfloor}\Psi_{n}^{r-1}(x-n-rq)2^{n-q}.\]
		Note that 		$q$ has its minimum value $0$ when  all $k_i\neq r$, and $q$ has  maximum value $\lceil\frac{x}{r+1}\rceil-1$ when all $k_i$'s except one, are equal to $r$.   Also note that there are ${n\choose q}$ possibilities of having $q$ number of $k_i$ to be equal to $r$ in the tuple $(k_1,k_2,\dots ,k_n)$. 
		Now summing up for $q$ in between $0$ and $\lceil\frac{x}{r+1}\rceil-1$, the total number of solutions is given by: 
		\[\sum_{q=0}^{\lceil\frac{x}{r+1}\rceil-1}\sum_{n=\lceil \frac{x-q}{r}\rceil}^{\lfloor \frac{x-(r-1)q}{2}\rfloor}\Psi_{n}^{r-1}(x-n-rq)2^{n-q}{n\choose q}.\]
		This proves the lemma. 
	\end{proof}
	
	
	\section{Counting the reciprocal conjugacy classes}\label{3}
	To count the reciprocal classes in the Hecke group, we need to specify the representations of the reciprocal classes. We recall the following theorem from \cite{DG},
	\begin{theorem}{\cite{DG}}\label{th}
		For any reciprocal hyperbolic element in $ \Gamma_{p} $, the following possibilities can occur.
		\begin{enumerate}
			\item If p is odd, every reciprocal class has exactly four symmetric elements.
			\item If p is even.
			\begin{enumerate}
				\item Suppose in the reciprocal class, the reciprocators are conjugate to $ \iota $. Then, there are exactly four symmetric elements in that class.
				\item Suppose in the reciprocal class, the reciprocators conjugate to $\gamma_{p}^{r} $. Then, there are exactly $ 2p $ many $p$-reciprocal elements. 
				\item Suppose in the reciprocal class, each reciprocal element has two types of reciprocators; one type of reciprocator conjugates to $ \iota $, and another type of reciprocator conjugates to $ \gamma_{p}^{r} $.
				\begin{enumerate}
					
					\item If the reciprocal class does not contain any non-zero power of $ \iota \g^r$, then the class has exactly two symmetric elements and $ p $ many $p$-reciprocal elements.
					
					
					\item If the reciprocal class contain any non-zero power of $ \iota \g^r$, then the class has exactly two symmetric p-reciprocal elements and $p-2$  $p$-reciprocal elements.			\end{enumerate}
			\end{enumerate}
		\end{enumerate}
		
	\end{theorem}
	
	\subsection{Counting reciprocal classes}
	We have seen that a reciprocal class $[\alpha]$ in $\Gamma_p$ can be divided into three types depending on whether it has all the reciprocators conjugate to only one of  $\iota$ and  $\gt$, or both of them. In the following, we obtain a `normal form' for such reciprocal classes. 
	Note that a reciprocal element $\alpha$ with the form  $ \io \gt\io \g^{k_{1}}\io \g^{k_{2}}\dots \io \g^{k_{n}}\io \gt\io \g^{-k_{n}}\dots\io \g^{-k_{2}}\io \g^{-k_{1}} $, is a symmetric $p$-reciprocal element, but any even power of $\alpha$ i.e., $\alpha^{2m}$ for  $m\in \z$, may not remain in the same form. It can take the form of a reduced expression like $\io k\io k^{-1}$ or  $\gt k \gt k^{-1}$, up to conjugacy. Although the reciprocal elements have distinct forms in the primitive case, they may coincide in the general case.
	\begin{lemma}\label{lem} Let $[{\tilde{\alpha}}]$ be a reciprocal class in $\Gamma_p$.  Then, we can choose a cyclically reduced representative $\alpha$ (or $\alpha^{-1}$) of this class as follows.  
		
		\begin{enumerate}
			\item Suppose all the reciprocators of $[\alpha]$ are conjugate to $\io$. Then  $\alpha$ is of the form  \\$ \io\g^{k_{1}}\io \g^{k_{2}}\dots \io \g^{k_{n}}\io \g^{-k_{n}}\dots\io \g^{-k_{2}}\io \g^{-k_{1}}  $.
			\item Suppose all the reciprocators of $[\alpha]$ are  conjugate to $ \gt $. Then $\alpha$ is of the form  \n  $ \io \gt\io \g^{k_{1}}\io \g^{k_{2}}\dots \io \g^{k_{n}}\io \gt\io \g^{-k_{n}}\dots\io \g^{-k_{2}}\io \g^{-k_{1}} $.
			\item Suppose  $[\alpha]$ has two types of reciprocators, some are conjugate to  $ \io $ and others conjugate to $ \gt $. 
			If the reciprocal class does not contain any power of $\io \gt$,   then $\alpha$ has the form  either \n
			$ \io\gt\io  \g^{k_{1}}\io \g^{k_{2}}\dots \io \g^{k_{n}}\io \g^{-k_{n}}\dots\io \g^{-k_{2}}\io \g^{-k_{1}} $ or $\io  \g^{k_{1}}\io \g^{k_{2}}\dots \io \g^{k_{n}}\io\gt\io \g^{-k_{n}}\dots\io \g^{-k_{2}}\io \g^{-k_{1}}$.
			
			
			If $[\alpha]$ contains any non-trivial power of $\io \gt  $, then it has a unique representative of the form $ (\io\gt )^{k}$ where $ k\in \z$. 
	\end{enumerate}\end{lemma}
	\begin{proof} Since the type of the class representative $\alpha$ remains unchanged under taking powers, except in the case where $\alpha$ is a symmetric $p$-reciprocal element, it suffices to prove the result for a primitive reciprocal class $[\alpha]$.
		
		For proving (1), assume that the reciprocal word $ R $ from the reciprocal class $  [\alpha]$  is a primitive symmetric element:  $R=\io \eta \io\eta^{-1}$, where $ \eta \in \Gamma_p $. There are exactly three other symmetric elements: $R^{-1}$,  
		$ \eta^{-1} R\eta$ and $\eta^{-1} R^{-1}\eta$. We can write $ \eta $, as a word in terms of $ \io \g $. Let, $ \eta_{1}= \io\g^{k_{1}}\io\g^{k_{2}}\dots \io\g^{k_{n}}$, then there are four choices of $ \eta $ either $ \eta_{1}, \io\eta_{1}, \eta_{1}\io \text{ or } \io \eta_{1}\io. $\n Our claim is that any  two of $R,~R^{-1}$,  
		$ \eta^{-1} R\eta$ and $\eta^{-1} R^{-1}\eta$ elements provide the required forms for the choices of $\eta$. 
		We first assume $ \eta=\eta_{1} $. Therefore,  \[R=\io\io\g^{k_{1}}\io\g^{k_{2}}\dots \io\g^{k_{n}}\io (\io\g^{k_{1}}\io\g^{k_{2}}\dots \io\g^{k_{n}})^{-1}\]
		\[=\g^{k_{1}}\io\g^{k_{2}}\dots \io\g^{k_{n}}\io\g^{-k_{n}}\io\dots\g^{-k_{2}}\io\g^{-k_{1}}\io.\]
		The above implies that $R^{-1}$ is in the desired form. 
		
		Also note that $\eta ^{-1}R^{-1}\eta$ is in similar form: $\io \g^{-k_{n}}\dots\io  \g^{-k_{2}}\io \g^{-k_{1}}\io\g^{k_{1}}\io \g^{k_{2}}\dots \io \g^{k_{n}}$. Thus, when  $ \eta=\eta_{1} $, the desired form of $\alpha$ is given by  either $ R^{-1}=\eta\io\eta^{-1}\io $ or $ \eta^{-1}R^{-1}\eta  = \io\eta^{-1}\io\eta $.
		
		Similarly, for the other choices of $\eta$, the desired form of $\alpha$ may be chosen as follows.
		
		When $\eta=\io \eta_1$, $\alpha$ is either $R$ or $\eta^{-1} R^{-1}\eta$; 
		\\
		for $ \eta=\eta_{1}\io $, $ R=\io(\eta_{1}\io)\io(\eta_{1}\io)^{-1}=\io \eta_{1}\io\eta_{1}^{-1}. $ So, similar to the first case, $\alpha$ has two choices $ R^{-1} $ and $ \eta^{-1}R\eta $.
		
		For $ \eta=\io\eta_{1}\io, $ $ R=\io\io\eta_{1}\io\io\io\eta_{1}^{-1}\io=\eta_{1}\io\eta_{1}^{-1}\io .$ So, similar to the second case, $\alpha$ have two choices, $ R $ and $ \eta^{-1}R\eta. $
		
		
		For $(2)$, assume that $R$ is a primitive  $p$-reciprocal element from $[\alpha].$ By \thmref {th}, there are $2p$ $p$-reciprocal elements, but we choose the cyclically reduced one as $R$. Let $R=\gt\eta \gt\eta^{-1}$.  Then, again by similar arguments as  in (1), if we choose, $ \eta= \io\g^{k_{1}}\io\g^{k_{2}}\dots \io\g^{k_{n}}$ then, 
		\[R=\gt \io\g^{k_{1}}\io\g^{k_{2}}\dots \io\g^{k_{n}}\gt \g^{-k_{n}}\io\g^{-k_{(n-1)}}\dots \io\g^{-k_{1}}\io\]
		\[\imp R=\gt \io\g^{k_{1}}\io\g^{k_{2}}\dots \io \g^{k_{n-1}}\io\gt\io\g^{-k_{n-1}}\dots \io\g^{-k_{1}}\io\]  
		\[\imp \io R\io =  \io \gt\io \g^{k_{1}}\io \g^{k_{2}}\dots \io \g^{k_{n}}\io \gt\io \g^{-k_{n}}\dots\io \g^{-k_{2}}\io \g^{-k_{1}}=R', \hbox{say}.\]	Thus  $R'$ is in the form of $\io \gt\io \g^{-k_{n}}\dots\io \g^{-k_{2}}\io \g^{-k_{1}} \io \gt\io \g^{k_{1}}\io \g^{k_{2}}\dots \io \g^{k_{n}}  $.  And $R''= \io\eta^{-1} R\eta \io$ is also in the same form. Since, $R$ and $\eta^{-1}R\eta$ are $ p $-reciprocal elements, in both cases $ \g^{k}R\g^{-k} $ and $ \g^{k} \eta^{-1}R\eta\g^{-k} $  for $-r<k_i\leq r$, are also $ p $-reciprocal elements that completes the list of $ p $-reciprocal elements which are not cyclically reduced. So, we have two choices of $\alpha$ from the cyclically reduced words:  $ R' , ~R'',R'^{-1}, R''^{-1}$. For other choices of $\eta$, as noted above,  the arguments will be similar as in the proof of $(1)$. This proves (2).
		
		For $(3)$, by \thmref{th}, there are $p+2$ choices for $R$. As even powers of 
		$R$ may already appear in the previous forms, we now focus on the remaining cases.  It is easy to see that only two of them are in the desired form:  $ R=\gt\eta\io \eta^{-1} $  and $ R'=\eta^{-1}\gt \eta \io.$ The  others are either $ \g^{k}R\g^{-k} $ or $\io R'\io$ where $-r<k\leq r$. Let $\eta_1$ as  in  $(1)$, so the choices of $\eta$ are, $\eta_1$, $\io\eta_1$, $\eta_1 \io$, and $\io\eta_1\io$.  
		If we assume $\eta=\eta_1$, then $\io R\io= \io\gt\eta\io\eta^{-1} $ and $R'=\eta^{-1}\gt\eta\io $ are the choices for $\alpha$. 
		By similar arguments as in  $ (1) $, we obtain the two desired forms for the remaining choices of $\eta$. This proves $(3)$.
	\end{proof}
	\begin{lemma}
		Consider the set of all reciprocal classes with length $2l$ and reciprocators conjugate to $\iota$. Then, this set has cardinality
		% $ \frac{1}{2}\sum_{n=\lceil\frac{l}{r+1}\rceil}^{\lfloor l/2 \rfloor}\Psi_{n} ^{r}(l-n)2^{n} $ if $ p $ is $2r+1$, otherwise
		$$\frac{1}{2}\sum_{q=0}^{\lceil\frac{l}{r+1}\rceil-1} \sum_{n= \lceil\frac{l-q}{r}\rceil }^{\lfloor \frac{l-(r-1)q}{2}\rfloor}\Psi_{n}^{r-1}(l-n-rq)2^{n-q} {n \choose q}.$$ 
	\end{lemma}
	\begin{proof}
		%	\tb{For the odd case:}
		%If $ p=2r+1 $ Every reciprocal class has a cyclically reduced word $ R= \io\g^{k_{1}}\io \g^{k_{2}}\dots \io \g^{k_{n}}\io \g^{-k_{n}}\dots\io \g^{-k_{2}}\io \g^{-k_{1}} $. Now we will calculate the number of reciprocal classes for each word length $ 2l $.
		%\[ \hbox{ length of word $ R $}=|R|=2\sum_{i=1}^{n} k_{i} +2n.\]
		%Let\begin{equation}\label{label}
			%	 2\sum_{i=1}^{n} k_{i} +2n=2l 
			%\end{equation}
			%\[\imp \sum_{i=1}^{n} k_{i} =l-n.\]
			%So, we just need to calculate the number for integral solutions for each $ k_{i}\leq r $.
			%Let us consider all $ k_{i} $'s  are non-zero then the number of solution \ref{label} is $ \Psi_{n} %^{r}(l-n) $(say). So, the number of all reciprocal classes of length $ 2l $ is $ \frac{1}{2}\sum_{n=\lceil\frac{l}{r+1}\rceil}^{\lfloor l/2 \rfloor}\Psi_{n} ^{r}(l-n)2^{n} $.
			%\tb{For the even case:}
			Consider the symmetric reciprocal word $ R= \io\g^{k_{1}}\io \g^{k_{2}}\dots \io \g^{k_{n}}\io \g^{-k_{n}}\dots\io \g^{-k_{2}}\io \g^{-k_{1}} $, where, $-r< k_{i}\leq r $ for all $1\leq i\leq n$. We see that 
			\[ \g^{-k_{n}}\io \dots \g^{-k_{1}}\io R\io\g^{k_{1}}\io \g^{k_{2}}\dots \io \g^{k_{n}}=\io \g^{-k_{n}}\dots\io \g^{-k_{2}}\io \g^{-k_{1}}\io\g^{k_{1}}\io \g^{k_{2}}\dots \io \g^{k_{n}} .\]
			Let  $R'=\io \g^{-k_{n}}\dots\io \g^{-k_{2}}\io \g^{-k_{1}}\io\g^{k_{1}}\io \g^{k_{2}}\dots \io \g^{k_{n}}$. Then 
			$ R $ and $ R' $ are of the form mentioned in \lemref{lem}. Since, given word length of $ R $ and  $ R' $ is $ 2l $, we have 
			\[\sum^n_{i=1} 2|k_i|+2n=2l ~\imp ~ \sum_{i=1}^{n} |k_{i}|=l-n, 		\hbox{ where } -r<k_{i}\leq r. \]
			Let $ q $ be the number of $ k_{i} $'s such that $k_i=r$. After rewriting the indices in the above equation, we get \[ \sum_{i=1}^{n-q} |k_{i}|=(l-n-rq). \]
			Then, by the similar arguments from \lemref{3.1},
			%\[\imp (l-n-rq)\leq (n-q)(r-1)\]
			%\[\imp l-n-rq\leq n(r-1)-qr+q\]
			%\imp l-q\leq nr\]
			%\[\imp \frac{l-q}{r}\leq n.\]
			the lower bound of $ n $ is $ \lceil\frac{l-q}{r}\rceil $ and 
			%When we add up all $ 1 $'s then this is the upper bound for $ n $, 
			%\[n-q=l-n-rq\]
			%\[\imp 2n=l-(r-1)q\]
			%\[\imp n=\frac{l-(r-1)q}{2}.\]
			the upper bound for $ n $ is $ \lfloor \frac{l-(r-1)q}{2}\rfloor $. Also, $ q $ varies from $ 0 $ to $\lceil\frac{l}{r+1}\rceil-1$, excluding the symmetric $ p $-reciprocal element $ (\io \gt)^{m} $ where, $ m\in \mathbb{Z} $. Now, using \lemref{lem}, \lemref{3.1} we obtain that the number of the reciprocal symmetric classes is  $$\frac{1}{2}\sum_{q=0}^{\lceil\frac{l}{r+1}\rceil-1} \sum_{n= \lceil\frac{l-q}{r}\rceil }^{\lfloor \frac{l-(r-1)q}{2}\rfloor}\Psi_{n}^{r-1}(l-n-rq)2^{n-q} {n \choose q}.$$
			This proves the assertion. 
		\end{proof}
		
		\medskip 
		%Note that, other type of reciprocal classes will be there only when $  p $ is even. So, to find the total number of reciprocal classes with a certain length, we also need to consider other cases. In the following lemma, we have found the number of conjugacy classes of $ p $-reciprocal elements.
		\begin{lemma}
			Consider the set of all  $p$-reciprocal classes with length $2l$. Then, this set has cardinality
			$$  \frac{1}{2}\sum_{q=0}^{(\lceil\frac{l}{r+1}\rceil-2)} \sum_{n=\lceil \frac{l-(r+1)-q}{r}\rceil}^{\lfloor\frac{l-1-(r+1)q-r}{2}\rfloor} \Psi_{n}^{r-1}(l-(n+1)-(q+1)r)2^{n-q} {n\choose q} ~ \hbox{ where }p=2r.$$
			
		\end{lemma}
		\begin{proof}
			Let $R$ be a $p$-reciprocal element of this form: $ \io \gt\io \g^{k_{1}}\io \g^{k_{2}}\dots \io \g^{k_{n}}\io \gt\io \g^{-k_{n}}\dots\io \g^{-k_{2}}\io \g^{-k_{1}}$, see, \lemref{lem}.
			Therefore,
			\[2\sum^{n-q}_{i=1}k_i+2n+2rq+2(r+1)=2l\]
			\begin{equation}
				\imp	\sum_{i=1}^{n-q} |k_{i}|+n+rq+ r+1=l, 
			\end{equation} where $ {|k_i|}\leq r-1$. Then, 
			\[ \sum_{i=-1}^{n-q}{|k_{i}|}+(n+1)+r(q+1)=l\]
			\[\imp \sum_{i=1}^{n-q} |k_{i}|=l-r(q+1)-(n+1).\]
			So, with the help of the last expression, we can deduce the total number of conjugacy classes. The remaining part is to show the upper and lower limits of $ n $ in the sum. 
			For the lower limit, 
			\[\sum_{i=1}^{n-q}|k_{i}|\leq (n-q)(r-1)\]
			\[\imp l-r(q+1)-(n+1) \leq (n-q)(r-1)\]
			\[\imp  l-r(q+1)-(n+1) \leq nr-n-qr+q \]
			\[\imp l-(r+1)-q\leq nr\]
			\[\imp \frac{l-(r+1)-q}{r}\leq n.\]
			So, the lower bound of $ n $ is $ \lceil \frac{l-(r+1)-q}{r}\rceil .$
			And the upper bound when all $ k_{i} $'s have the lowest value, which is $ 1 $, i.e.,
			\[n-q=l-(n+1)-r(q+1)\]
			\[\imp n=\frac{l-1-(r+1)q-r}{2}.\]
			So, the upper bound of $n$ is $ \lfloor\frac{l-1-(r+1)q-r}{2}\rfloor. $ And, $q$ varies from $0$ to $\lceil\frac{l-r-1}{r+1}\rceil-1=\lceil\frac{l}{r+1}\rceil-2$.  Then by \lemref{3.1},\lemref{lem}, the number of $ p $-reciprocal classes where the reciprocators are conjugate of $\gt$, is 
			$$  \frac{1}{2}\sum_{q=0}^{(\lceil\frac{l}{r+1}\rceil-2)} \sum_{\lceil \frac{l-(r+1)-q}{r}\rceil}^{\lfloor\frac{l-1-(r+1)q-r}{2}\rfloor} \Psi_{n}^{r-1}(l-(n+1)-(q+1)r)2^{n-q} {n\choose q}.$$
			This proves the assertion. 
		\end{proof}
		\begin{lemma}\label{lem1}
			The number of symmetric $p$-reciprocal classes that does not contain any non-trivial power of $ \io\gt $ of word length $ 2l $, resp.  $2l+1$,  is $$\frac{1}{2}\sum_{q=0}^{\lceil\frac{l-u}{r+1}\rceil-1} \sum_{n= \lceil\frac{l-u-q}{r}\rceil }^{\lfloor \frac{l-u-(r-1)q}{2}\rfloor}\Psi_{n}^{r-1}(l-u-n-rq)2^{n-q} {n \choose q} $$ where, $r=2u $, resp. $r=2u+1$.
		\end{lemma}
		\begin{proof}
			Suppose $R$ is the cyclically reduced representation of the mentioned reciprocal class, then \begin{equation} R= \io \gt\io \g^{k_{1}}\io \g^{k_{2}}\dots \io \g^{k_{n}}\io \g^{-k_{n}}\dots\io \g^{-k_{2}}\io \g^{-k_{1}}.\end{equation} 
			Here, we first consider  $r=2u$. Then
			\begin{equation}
				r+1+ 2\sum_{i=1}^{n} |k_{i}| +2n=2l+1
			\end{equation} where $k_{i}\leq r-1$ for all $i$. Then implies,
			\[	 2\sum_{i=1}^{n} |k_{i}| +2n=2(l-u) \]
			\[\imp 	 \sum_{i=1}^{n} |k_{i}| +n=l-u. \]
			Also, $R$ is conjugate to $\io  \g^{k_{1}}\io \g^{k_{2}}\dots \io \g^{k_{n}}\io\gt\io \g^{-k_{n}}\dots\io \g^{-k_{2}}\io \g^{-k_{1}}$, so there are two choices of words for the required form.
			Now, the process is the same as the symmetric element case. So, the number of the reciprocal class is $\frac{1}{2}\sum_{q=0}^{\lceil\frac{l-u}{r+1}\rceil-1} \sum_{n= \lceil\frac{l-u-q}{r}\rceil }^{\lfloor \frac{l-u-(r-1)q}{2}\rfloor}\Psi_{n}^{r-1}(l-u-n-rq)2^{n-q} {n \choose q} $.
			
			
			Similarly, for the odd case, if we choose $u=\lfloor \frac{r-1}{2}\rfloor$, we will get same number of conjugacy classes. 		\end{proof}
		
		
		\medskip 		The only remaining case when the reciprocal class is represented by a non-trivial power of $\io\gt $, follow from \lemref{lem}. Combining all these, we have the following theorem. 
		\begin{thm}\label{count1}
			Let $\Gamma_{p}$ be the Hecke group with the presentation $\langle \io,\g~|~\io^{2},\g^{p}\rangle$ and $p=2r.$ Then the following statements hold:
			\begin{enumerate}
				\item The number of symmertric reciprocal classes of word length $2l$ is, 	$$\frac{1}{2}\sum_{q=0}^{\lceil\frac{l}{r+1}\rceil-1} \sum_{n= \lceil\frac{l-q}{r}\rceil }^{\lfloor \frac{l-(r-1)q}{2}\rfloor}\Psi_{n}^{r-1}(l-n-rq)2^{n-q} {n \choose q}.$$ 
				\item The number of $p$-reciprocal classes of word length $2l$ is,
				$$  \frac{1}{2}\sum_{q=0}^{(\lceil\frac{l}{r+1}\rceil-2)} \sum_{n=\lceil \frac{l-(r+1)-q}{r}\rceil}^{\lfloor\frac{l-1-(r+1)q-r}{2}\rfloor} \Psi_{n}^{r-1}(l-(n+1)-(q+1)r)2^{n-q} {n\choose q} .$$
				
				\item The number of symmetric $ p$-reciprocal classes of word length $2l$, resp. $2l+1$ depending on $r$ is,
				$$\frac{1}{2}\sum_{q=0}^{\lceil\frac{l-u}{r+1}\rceil-1} \sum_{n= \lceil\frac{l-u-q}{r}\rceil }^{\lfloor \frac{l-u-(r-1)q}{2}\rfloor}\Psi_{n}^{r-1}(l-u-n-rq)2^{n-q} {n \choose q} +1$$ where, $r=2u $, resp. $r=2u+1.$
			\end{enumerate}
		\end{thm}
		\begin{proof}
			The proof of this proposition follows from the previous lemmas.
		\end{proof}
		\section{Spherical growth  of the Reciprocal Classes}\label{4}
		In this section, we prove our main result. 
		\subsection{When $r$ is odd }
		In this case, $r=2u+1$ when the reciprocal classes always have even word lengths, since the words representing the reciprocal classes all have even word length. 
		
		\begin{lemma}\label{count}
			Let $\mathcal{N}_{2l}$ be the set of reciprocal classes of word length $2l$, then the following condition will hold.
			%$\io\g^{k_{1}}\io\g^{k_{2}}\dots \io\g^{k_{n}}\io \g^{-k_{n}}\dots \io\g^{-k_{1}}$.
			\begin{enumerate}
				\item If $l\leq r$, then the number of reciprocal classes is $ \frac{1}{6}(2^{l}+2(-1)^{l})$.
				\item If $l=r+1$, the number of reciprocal classes is $\frac{1}{6}(2^{l}+2(-1)^{l})+\frac{ 1}{6}(2^{u+1}+2(-1)^{u+1})-1.$
				\item If $l\geq r+2$, 
				\[|\mathcal{N}_{2l}|=2\sum_{j=1}^{r-1}|\mathcal{N}_{2(l-j-1)}|+|\mathcal{N}_{2(l-r-1)}|.\]
			\end{enumerate} 
		\end{lemma}
		\begin{proof}
			For $(1)$, it is given that  $l\leq r$.  There is only a symmetric reciprocal class; otherwise, the length will exceed $2r$.
			So, the conditions to determine the cardinality of $\mathcal{N}_{2l}$ for $\z_{2}*\z_{p}$ is same as the cardinality of $\mathcal{N}_{2l}$ for $\z_{2}*\z_{p-1}$. Since, $p-1$ is odd, then from \cite[Theorem 5.3.2.]{Ma}, the cardinality of the cyclically reduced reciprocal words is 
			\[\frac{1}{3}(2^{l}+2(-1)^{l})\]
			Since each reciprocal class has two cyclically reduced reciprocal words, this proves $(1)$.			
			
			For $(2)$,
			again, there is no $p$-reciprocal class in $\mathcal{N}_{2l}$ otherwise the word length will exceed $2(r+1).$ So, the conjugacy classes are either symmetric reciprocal classes or primitive symmetric $p $-reciprocal classes. Let $R$ be the representative of a conjugacy class in $\mathcal{N}_{2l}$ as mentioned in the proof of the \ref{lem}. Then, \[R=(\io\gt)^{\epsilon}\io  \g^{k_{1}}\io \g^{k_{2}}\dots \io \g^{k_{n}}\io \g^{-k_{n}}\dots\io \g^{-k_{2}}\io \g^{-k_{1}} \] where, $\epsilon$ is $0$ or $1$ on the respective choice of symmetric reciprocal class or symmetric $ p$-reciprocal class.  In any case, we get this equation for the word length $2l$,
			\begin{equation}
				2\sum_{i=1}^{n}|k_{i}| +2n+\epsilon(r+1)= 2l.
			\end{equation}
			Let us consider $q$ as the number of $ k_i$'s such that they are equal to $r.$ So,  after rewriting the $k_{i}$'s, we have,
			\begin{equation}
				2\sum_{i=1}^{n}|k_{i}| +2n+2qr+\epsilon(r+1)= 2l
			\end{equation} where, $k_{i}\leq r-1$ and $q\geq 0$, $n>0$ and $r=2u+1. $ 
			This implies, \[2\sum_{i=1}^{n-q}|k_{i}| +2n+2qr+\epsilon(r+1)= 2(r+1).\]
			Three disjoint cases may arise.
			
			\medskip  \tb{Case 1:}	If $\epsilon=0$ and $q=0$,
			then, the element will belong to symmetric classes in $\mathcal{N}_{2l}.$ In that case, we can get $|\mathcal{N}_{2l}|$ from \cite[Theorem 5.3.2.]{Ma} i.e., $\frac{ 1}{3}(2^{l}+2(-1)^{l})-2$ using same fact as $(1)$. 
			
			\tb{Case 2:} 	If	$\epsilon=1$ and $q=0$, this implies, \[2\sum_{i=1}^{n}|k_{i}| +2n+(r+1)= 2(r+1).\]
			\[\imp 2\sum_{i=1}^{n}|k_{i}| +2n+=(r+1) \]
			Therefore, $\sum_{i=1}^{n}|k_{i}| +n=(u+1)$, 
			where $k_{i}\leq r-1.$ Since, $u+1\leq r$ then, from $(1)$ we obtain the value, $\frac{ 1}{3}(2^{u+1}+2(-1)^{u+1})$,  where $u=\frac{r-1}{2}.$
			
			
			\tb{Case 3:} If $q=1$ and $\epsilon=1,$ this implies
			\[ 2\sum_{i=1}^{n-1}|k_{i}| +2n+2r= 2(r+1)\]
			\[\imp\sum_{i=1}^{n-1}|k_{i}| +n+r= (r+1)\]
			\[\imp \sum_{i=1}^{n}|k_{i}|=(1-n). \]
			Since absolute values are non-negative, then $n$ and $q$ must be $1$, and all $k_{i}$'s are $0.$ Then $R=\io\gt\io\gt$. So, there is only one element for the case $q>0.$		
			Then the required number is $\frac{ 1}{3}(2^{l}+2(-1)^{l})+\frac{ 1}{3}(2^{u+1}+2(-1)^{u+1})-1.$
			It follows the proof of (2).
			
			For(3), it suffices to consider primitive case since all arguments are based on the powers of $ \g .$ We divide $\mathcal{N}_{2l}$ into $3$ disjoint sets such as $A_{2l},B_{2l},C_{2l}$ such that,
			$A_{2l}, ~B_{2l},~C_{2l}$ represent the set of symmetric reciprocal classes, $p$-reciprocal classes and symmetric $p$-reciprocal classes respectively. We can see that the set $A_{2l}$ is a collection of the reduced words whose primitive elements are symmetric of the form  $  \io\g^{k_{1}}\io \g^{k_{2}}\dots \io \g^{k_{n}}\io \g^{-k_{n}}\dots\io \g^{-k_{2}}\io \g^{-k_{1}} $ where, $-r<k_{i}\leq r$. Then, the set $A_{2l,j}=\io\g^{j}\io \g^{k_{2}}\dots \io \g^{k_{n}}\io \g^{-k_{n}}\dots\io \g^{-k_{2}}\io \g^{-j}$  gives the equation,
			\[ 2+2|j|+2\sum_{i=2}^{n}|k_{i}|+2n=2l\]
			where, $-r<k_i,~ j\leq r$ .
			This implies, \[\sum_{i=2}^{n}|k_{i}|+n=l-|j|-1.\]
			So, $|A_{2l,j}|=|A_{2(l-|j|-1)}|$. Now, varying  and $|j|$ over varies in the set $ \{r,r-1,\dots, 1\} $, then we have this following condition will hold,
			\[|A_{2l}|=2\sum_{j=1}^{r-1}|A_{2(l-|j|-1)}|+|A_{l-r-1}|.\]
			Here, $ 2 $ appears for both signs of $k_{i}$'s.
			Similarly, we have the following equalities.
			$$|B_{2l}|=2\sum_{j=1}^{r-1}|B_{2(l-j-1)}|+|B_{l-r-1}|, ~~ |C_{2l}|=2\sum_{j=1}^{r-1}|C_{2(l-j-1)}|+|C_{l-r-1}|.$$ 
			Therefore, we have
			\[|\mathcal{N}_{2l}|=2\sum_{j=1}^{r-1}|\mathcal{N}_{2(l-j-1)}|+|\mathcal{N}_{l-r-1}|.\]
			This proves the lemma. 
		\end{proof}
		
		\medskip The characteristic equation of the given recurrence relation, 
		$|\mathcal{N}_{2l}|=2\sum_{j=1}^{r-1}|\mathcal{N}_{2(l-j-1)}|+|\mathcal{N}_{2(l-r-1)}|$ is
		\begin{equation}
			x^{l}=2\sum_{j=1}^{r-1}x^{l-j-1}+x^{l-r-1}.
		\end{equation} The non-zero solutions from the above equations, will also satisfy the following equation:
		\begin{equation} \label{re} x^{r+1}-2\sum_{j=1}^{r-1}x^{r-j}-1=0.\end{equation} 
		
		
		By Descarte's rule of signs, there is only one real positive root. Now, we use the following theorems by Gauss to investigate the roots of \eqnref{re}. 
		\begin{theorem}\label{4.2}
			Every real root of a monic polynomial with integer coefficients is either an integer or irrational.
		\end{theorem}
		
		\begin{defn}
			Let $f(z)=z^n+a_{n-1}z^{n-1}+a_{n-2}z^{n-2}+\dots +a_1 z+a_0$ be an monic  with complex coefficients. Then, the \textit{Cauchy bound} of $f$, say $r$, is the unique positive root of the mentioned polynomial. 
		\end{defn}
		\begin{thm}\label{4.4}
			Let  $f(z)=z^n+a_{n-1}z^{n-1}+a_{n-2}z^{n-2}+\dots +a_1 z+a_0$  be an $n$-degree polynomial with complex coefficients. Then all zeros of $f(z)$ have modulus less than or equal to $r.$
		\end{thm}
		
		
		
		\begin{lemma}\label{2.7}
			The following recurrence relation   \begin{equation}\label{2.7}
				|\mathcal{N}_{2l}|=2\sum_{j=1}^{r-1}|\mathcal{N}_{2(l-j-1)}|+|\mathcal{N}_{2(l-r-1)}|\end{equation}  has a convergent limit in between $\sqrt{2}$ and $2$.
			
			Moreover, $|\mathcal{N}_{2l}|=\mathcal{O}( l^{s-1}\rho^{l})$ as $l\rightarrow \infty$, where $s$ is the maximum of the multiplicities of roots of the recurrence polynomial.
		\end{lemma}	
		
		\begin{proof}
			The recurrence polynomial $p(x)= x^{r+1}-2\sum_{j=1}^{r-1}x^{r-j}-1$ of the relation \ref{2.7}. If we define $p$ as a map from real line to real line, so $p(\sqrt{2})<0$: 
			\[ p(\sqrt{2})=2^{\frac{r+1}{2}}-2\sum_{j=1}^{r-1}2^{\frac{r-j}{2}}-1\]
			\[=2^{\frac{r+1}{2}}-2^{\frac{r+1}{2}}-2\sum_{j=2}^{r-1}2^{\frac{r-j}{2}}-1\]
			\[=-2\sum_{j=2}^{r-1}2^{\frac{r-j}{2}}-1<0.\] 		
            Also, 
			\[p(2)= 2^{r+1}-2\sum_{j=1}^{r-1}2^{r-j}-1\]
			\[=2^{r+1}-2(2^{r-1}+2^{r-2}+\dots +2^{2}+2+1)+1\]
			\[=2^{r+1}-2(2^{r}-1)+1=3.\]
			That means, $p(2)>0.$
			
			Then, by Bolzano's theorem, there is at least one zero of the polynomial in the interval $[\sqrt{2},2]$, denote it as $\rho.$  Since, there is no integer in between $[\sqrt{2},2]$, $\rho$ must be an irrational number by \lemref{4.2}. This is the Cauchy bound for $p(x)$. Let the roots of the equation $p(x)=0$ are $ \rho_{1},\rho_{2}\dots \rho_{m} $ where, $ m\leq r+1 $.
			Then, the closed form to the recurrence relation is,
			\[|\mathcal{N}_{2l}|=c_{1}\rho_{1}^{l}+c_{2}\rho_{2}^{l}+\dots +c_{m}\rho_{m}^{l}. \]
			$\rho$ is the root of $p(x)=0$ with the largest absolute value by \thmref{4.4}. And $ c_{k} $ be the coefficient of the repeated root with the largest multiplicity,  say $s$. It is well known that $c_{k}$ is a polynomial in $l$ of degree $(s-1)$. Then, $\dfrac{|\mathcal{N}_{2l}|}{l^{s-1}\rho^{l}}\rightarrow finite$ as $l\rightarrow \infty.$ Therefore,  ${|\mathcal{N}_{2l}|}=\mathcal{O}({l^{s-1}\rho^{l}})$ as $l\rightarrow\infty$.
		\end{proof}
		%\begin{lemma}\label{2.8}
		%	$R^{np}_{2l}$ is the set of the non-primitive reciprocal classes, then $|R^{np}_{2l}|=\mathcal{O} \gamma_{1}c_{k}l\rho^{l/2}$
		%\end{lemma}
		\subsection{When $r$ is even }
		Odd word length reciprocal geodesic occurs when $r= 2u$. So, we will consider this condition throughout this subsection.
		
		\begin{lemma}
			Suppose $ \mathcal{N}_{2l-1}$ be the set of conjugacy classes of reciprocal words of word length $2l-1$. Then we have :
			
			\begin{enumerate}
				\item If $l\leq r+u+1$, then, $|\mathcal{N}_{2l-1}|=\frac{1}{6}(2^{l-u-1}+2(-1)^{l-u-1})$.
				\item If $l= r+u+2$, then $|\mathcal{N}_{2l-1}|=\frac{1}{6}(2^{l-u-1}+2(-1)^{l-u-1})+\frac{ 1}{6}(2^{u+1}+2(-1)^{u+1})-1$.
				\item If $l\geq r+u+3$, 
				\[|\mathcal{N}_{2l-1}|=|\mathcal{N}_{2(l-u-1)}|=2\sum_{j=1}^{r-1}|\mathcal{N}_{2(l-u-j-2)}|+|\mathcal{N}_{2(l-u-r-2)}|.\]
			\end{enumerate}
		\end{lemma}
		\begin{proof}
			To prove $ (1) $ it suffices to show that $|\mathcal{N}_{2l-1}|=|\mathcal{N}_{2(l-u-1)}|$. Let us consider the reciprocal reduced word with word length $2l-1$,  $$ (\io\gt\io  \g^{k_{1}}\io \g^{k_{2}}\dots \io \g^{k_{n}}\io \g^{-k_{n}}\dots\io \g^{-k_{2}}\io \g^{-k_{1}})^\epsilon $$ where, $ \epsilon $ is always odd natural number. But here, $\epsilon$ must be $ 1 $; otherwise, the word length will exceed $ 3r $, which will be a contradiction. Then, we get this equation for the word length $2l-1$ after rewriting the $k_{i}$'s,
			\[ 2\sum_{i=1}^{n}|k_{i}| +2n+2qr+2u+1= 2l-1\]
			\[\imp\sum_{i=1}^{n}|k_{i}| +n+qr= l-1-u. \]
			$ (1)$ follows from the above equality. 
			
			
			$ (2) $ follows from the previous case and from the arguments in the proof of \lemref{count} and the fact that $ [(\io\gt)^{3}] $ also belongs to $ \mathcal{N}_{3(r+1)}. $
			
			For $ (3) $, consider the primitive case; then, this will follow from the proof of \lemref{count}. \end{proof}
		
		\medskip 
		\begin{pro}
			A non-constant polynomial $p(x)$ in $\mathbb{Z}[X]$ is irreducible if and only if $p(x+a)$ is irreducible in $\mathbb{Z}[X]$ where $a$ is an integer. 
		\end{pro}
		So, $$p(x+1)=(x+1)^{r+1}-2\sum_{j=1}^{r-1}(x+1)^{r-j}-1$$
		The constant term of the polynomial is, $1-2\underbrace{(1+\dots+1)}_{\text{odd times}}-1$, say $2c$ where, $c$ is odd.
		Then, by the above lemma and applying Eisenstein Criterion with the prime $2$, we obtain that $p(z)$ is irreducible in $\mathbb{Z}[X]$. Now note the following well-known theorem: 
		
		\begin{thm}
			A non-constant polynomial in $ \mathbb{Z}[\textit{X}]$ is irreducible in $\mathbb{Z}[\textit{X}]$ if and only if it is both irreducible in $\mathbb{Q}[\textit{X}]$ and primitive in $\mathbb{Z}[\textit{X}]$ .
		\end{thm}
		It follows that $p(z)$ is irreducible in $\mathbb{Q}[X].$ It is also well-known that an irreducible polynomial over $\mathbb{Q}$ cannot have repeated roots in $\mathbb{C}$.  
		Therefore, $p(z)$ has no repeated roots in $\mathbb{C}$. This implies all $c_k$'s are constant, i.e. $s$ must be $1$. 
		
		\begin{lemma}\label{5} $|\mathcal{N}_{2l-1}|=\mathcal{O} (\rho^{l})$ as $l\rightarrow \infty$ where, $ \rho $ is the real positive root of the equation $p(z)=0$. 
		\end{lemma}
		\begin{proof}
			By the above arguments and  \lemref{2.7}, as $l\rightarrow \infty$ we see that $$|\mathcal{N}_{2l-1}|=\mathcal{O} (l^{s-1}\rho^{l-u-1})=\mathcal{O}(\rho^{l}),$$ since, $|\rho^{u}|$ is constant and $s$ is $1$.
		\end{proof}
		
		\subsection*{Proof of \thmref{mthm}} 
		The theorem follows from \lemref{2.7} and \lemref{5}. \qed
		\medskip

To establish \corref{cor1}, we observe the following lemmas. 
	\begin{lemma}
		Let $j = \lfloor \tfrac{l}{2} \rfloor$. Then
		\[
		|\mathcal{N}^{np}_{l}| = \mathcal{O}\left(\left\lfloor \tfrac{j+1}{2} \right\rfloor^{s-1} \rho^{\left\lfloor \tfrac{j+1}{2} \right\rfloor} \right),
		\]
		where $\mathcal{N}^{np}_{l}$ denotes the set of all non-primitive reciprocal classes in $\Gamma_p$ with word length $l$.
	\end{lemma}
	
	\begin{proof}
		Let $\mathcal{N}^{p}_{l}$ denote the set of all primitive reciprocal classes in $\Gamma_p$ with word length $l$. Clearly, both $\mathcal{N}^{np}_{l}$ and $\mathcal{N}^{p}_{l}$ are subcollections of $\mathcal{N}_{l}$. Note that a non-primitive reciprocal class of word length $l$ arises from a primitive class of word length $q$ dividing $l$, where $q < l$. Therefore,
		\[
		|\mathcal{N}^{np}_{l}| \leq \sum_{q \mid l,~q < l} |\mathcal{N}^{p}_{q}| \leq \sum_{q = 1}^{\lfloor \frac{l}{2} \rfloor} |\mathcal{N}_{q}|.
		\]
		Using the known upper bound on $|\mathcal{N}_q|$, we obtain
		\[
		\sum_{q=1}^{\lfloor \frac{l}{2} \rfloor} |\mathcal{N}_{q}| = \sum_{q=1}^{\lfloor \frac{l}{2} \rfloor} \mathcal{O}\left(\left\lfloor \tfrac{q+1}{2} \right\rfloor^{s-1} \rho^{\left\lfloor \tfrac{q+1}{2} \right\rfloor} \right).
		\]
		Since the dominant term in this sum corresponds to the largest value of $q$, namely $q = \lfloor \tfrac{l}{2} \rfloor = j$, we conclude that
		\[
		|\mathcal{N}^{np}_{l}| = \mathcal{O}\left(\left\lfloor \tfrac{j+1}{2} \right\rfloor^{s-1} \rho^{\left\lfloor \tfrac{j+1}{2} \right\rfloor} \right).
		\]
		This proves the lemma. 
	\end{proof}
	
		
		
		
		
		
		
		
		% number of reciprocal classes (whether it is odd or even) is asymptotic to an exponential function. And the primitive and non-primitive results are also similar but not equivalent.\\
		%	Note that here, the reciprocal words have with even and odd word length both. So, the counting of non-primitive reciprocal elements with even word length is not same result as \lemref{2.8} since, the non-primitive elements of word length $2l$ can be seen as power of an primitive element of word length even. 
		%\begin{lemma}
		%	$R^{np}_{2l-1}$ be the collection of reciprocal classes with word length $2l-1$. Then, $R^{np}_{2l-1}\sim \gamma_{1}c_{k}(l-u-1)\rho^{l/2}$.
		
		%\end{lemma}
		%\begin{proof}
		%	The proof is similar to \lemref{2.8}.
		%\end{proof}	\section{new}
	\begin{lemma}\label{prim}
			Let $\mathcal{N}_{l}$ be the set of all reciprocal classes in $\Gamma_p$ with word length $l$, and $\mathcal{N}^{p}_{l}$ be the subset of primitive reciprocal classes. Then \[
		|\mathcal{N}^{p}_{l}| \sim |\mathcal{N}_{l}|.
		\]
		
	\end{lemma}
	
	\begin{proof}
		Let $m$ be the multiplicity of the dominant root $\rho$ of the recurrence relation defined by the characteristic equation \[x^{r+1} - 2\sum_{j=1}^{r-1} x^{r-j} - 1 = 0.\]Then the total number of reciprocal classes satisfies
		\[
		|\mathcal{N}_{l}| \sim \left(\left\lfloor \tfrac{l+1}{2} \right\rfloor^{m} \rho^{\left\lfloor \tfrac{l+1}{2} \right\rfloor} \right),
		\]
		and the number of non-primitive classes satisfies the upper bound
		\[
		|\mathcal{N}^{np}_{l}| \leq \left(\left\lfloor \tfrac{j+1}{2} \right\rfloor^{m} \rho^{\left\lfloor \tfrac{j+1}{2} \right\rfloor} \right)
		\quad \text{for } j = \left\lfloor \tfrac{l}{2} \right\rfloor.
		\]
	
		Also, by definition,
		\[
		|\mathcal{N}_{l}| = |\mathcal{N}^{np}_{l}| + |\mathcal{N}^{p}_{l}| \leq \left(\left\lfloor \tfrac{j+1}{2} \right\rfloor^{m} \rho^{\left\lfloor \tfrac{j+1}{2} \right\rfloor} \right) + |\mathcal{N}^{p}_{l}|.
		\]
		We obtain the inequality
		\[
		|\mathcal{N}_{l}| - \left(\left\lfloor \tfrac{j+1}{2} \right\rfloor^{m} \rho^{\left\lfloor \tfrac{j+1}{2} \right\rfloor} \right) \leq |\mathcal{N}^{p}_{l}| \leq |\mathcal{N}_{l}|.
		\]
		Dividing by $|\mathcal{N}_{l}|$, we get
		\[
		1 - \frac{\left(\left\lfloor \tfrac{j+1}{2} \right\rfloor^{m} \rho^{\left\lfloor \tfrac{j+1}{2} \right\rfloor} \right)}{|\mathcal{N}_{l}|}
		\leq \frac{|\mathcal{N}^{p}_{l}|}{|\mathcal{N}_{l}|} \leq 1.
		\]
		
		Since \( j = \lfloor \tfrac{l}{2} \rfloor \) and both the numerator and denominator grow like a power of $\rho$ with comparable exponents as \( l \to \infty \), the ratio tends to zero. Thus,
		\[
		\frac{|\mathcal{N}^{p}_{l}|}{|\mathcal{N}_{l}|} \to 1 \quad \text{as } l \to \infty,
		\]
		which implies
		\[
		|\mathcal{N}^{p}_{l}| \sim |\mathcal{N}_{l}|.
		\]
    This proves the lemma. 
	\end{proof}

\subsection{Proof of \corref{cor1}}
	The corollary follows from \lemref{prim} and \thmref{mthm}. \qed
		
		
		\subsection*{Acknowledgment} 
		Das acknowledges support from UGC-NFSC senior research fellowship. Gongopadhyay is partially supported by the SERB core research grant CRG/2022/003680.
		
		The authors thank Anuj Jakhar for useful discussions on the roots of polynomials over the integers. 
		
		\begin{thebibliography}{99}
			\bibitem[BR]{BR} M. Baake, J. A. G. Roberts,
			\newblock The structure of reversing symmetry groups. 
			\newblock Bull. Aust. Math. Soc.  73, No. 3 (2006),  445-459. 
			\bibitem[BV]{BV}A. Basmajian and R. Suzzi Valli, \emph{Combinatorial Growth in the Modular Group},
			\newblock {Groups Geom. Dyn. 16 (2022), 683–703}. 
			
			\bibitem[BK]{BK} J. Bourgain and A.  Kontorovich, 
			\emph{Beyond expansion, III: Reciprocal geodesics}.
			Duke Math. J.168 (2019), no.18, 3413–3435.
			
			\bibitem[DG]{DG} D. Das and  K. Gongopadhyay.
			\newblock{\emph{Reciprocity in the Hecke Groups}.}
			\newblock     New York J. Math. 30 (2024) 722–744.
			
			\bibitem[ES]{ES} V. Erlandsson and J. Souto,{\em Counting and equidistribution of reciprocal geodesics and dihedral groups.} \newblock Int. Math. Res. Not.  (2024),  no. 13, 10298–10318.
			
			\bibitem[LR]{LR}J. S. W. Lamb and J. A. G. Roberts,{\em Time-reversal symmetry in dynamical systems: a survey}. Time-reversal symmetry in dynamical systems (Coventry, 1996). Phys. D 112:1–2 (1998), 1–39. 
			
			
			\bibitem[Ma]{Ma} B. T. Marmolejo, 
			\newblock {\emph{Growth of Conjugacy Classes of Reciprocal Words in Triangle Groups}(2020).} Ph.D. Thesis, City University of New York. 
			\newblock CUNY Academic Works.
			\newblock \url{https://academicworks.cuny.edu/gc_etds/3972}
			
			
			
			\bibitem[OS]{OS} A.  G. O’Farrell and I.  Short, {\it Reversibility in Dynamics and Group Theory}. London Mathematical Society Lecture Note Series, vol.416, Cambridge University Press, Cambridge, 2015.

            \bibitem[PP]{PP}J. Parkkonen and F. Paulin, {\em Counting and equidistribution of strongly reversible closed geodesics in negative curvature}. arXiv:2505.07738. 
			
			\bibitem[Sa]{Sa} P. Sarnak, 
			\newblock {\em{Reciprocal geodesics}.} 
			\newblock Analytic number theory,  
			Clay Math. Proc. 7,  217–237, Amer. Math. Soc., Providence, RI, 2007. 
			
			
			
			
			
			
			
		\end{thebibliography}
	\end{document}
